\section{Preliminaries}
\label{sec:prelim}
In this section, we gather some background on Bohr sets, kernels and homomorphisms in compact abelian groups. Most of the results are well-known or resemble known theorems. We include proofs for the results that we cannot pinpoint precisely in the literature. 

% We prove the results whose proofs we cannot find in literature.

\subsection{Notation} We write $[N]$ for the set $\{ 1, \ldots, N\}$. If $A$ and $B$ are two quantities, we write $A = O(B)$ or $A \ll B$ if there is a constant $C$ such that $|A| \leq CB$. We write $e(x)$ for $e^{2 \pi i x}$.

Throughout this paper, $G$ is a Hausdorff compact abelian group with probability Haar measure $\mu$ and  $\Gamma$ is the dual of $G$, written additively. The relevance of homomorphisms is that if $\gamma \in \Gamma$ and $\phi:G \rightarrow G$ is a continuous homomorphism, then $\gamma \circ \phi$ is also an element of $\Gamma$.

If $f: G \rightarrow \C$ is a function, for $t \in G$ we define the function $f_t(x) = f(x+t)$. 
For $f \in L^1(G)$, the Fourier transform of $f$ is the function
\[
\widehat{f} (\gamma) = \int_{G} f(x) \overline{\gamma (x)} \, d\mu(x) \qquad \textup{ for } \gamma \in \Gamma.
\]
For $f, g \in L^2(G)$, we then have Parseval's formula
\[
\int_{G} f(x) \overline{g(x)} \, d\mu(x) = \sum_{\gamma \in \Gamma} \widehat{f} (\gamma) \overline{ \widehat{g} (\gamma) }
\]
and Parseval's formula
\[
\int_{G} |f(x)|^2 \, d\mu(x) = \sum_{\gamma \in \Gamma} \left| \widehat{f} (\gamma) \right|^2.
\]


\subsection{Bohr sets} 
% If $\Lambda \subset \Gamma$ is a finite set and $\eta >0$, then the set
% \[
% B(\Lambda; \eta) = \{ x \in G: |\gamma(x) -1 | < \eta \textup{ for all } \gamma \in \Lambda \}
% \]
% is called a Bohr-$(|\Lambda|, \eta)$ set.
For $\Lambda, \Lambda_1, \Lambda_2 \subseteq \Gamma$ and $\eta_1, \eta_2 > 0$, it follows from the definition of Bohr sets that 
\[
    B(\Lambda_1; \eta_1) \cap B(\Lambda_2; \eta_2) \supset B(\Lambda_1 \cup \Lambda_2; \min(\eta_1, \eta_2))
\]
and
\[
    B(\Lambda; \eta_1) + B(\Lambda; \eta_2) \subset B(\Lambda; \eta_1 + \eta_2).
\]

\begin{lemma} \label{lem:translate}
Suppose $f_1, \ldots, f_k \in L^\infty(G)$, $\| f_i \|_{\infty} \leq 1$ for all $i=1, \ldots, k$. Let $\phi_1, \ldots, \phi_k$ be continuous homomorphisms $G \rightarrow G$. Then for any $\eta >0$, the set
\[
B = \{ t \in G: \| \widehat{f_i} - \widehat{f_{i, \phi_i(t)}} \|_{\infty} < \eta \textup{ for }i=1, \ldots, k \}
\]
contains a Bohr set $B(\Lambda; \eta)$ where $|\Lambda| \leq \frac{4k}{\eta^2}$.
\end{lemma}
\begin{proof}
Note that if $\| \widehat{f_i} - \widehat{f_{i, \phi_i(t)}} \|_\infty \geq \eta$, then for some $\gamma \in \Gamma$, 
\[
|\widehat{f_i}(\gamma) - \widehat{f_{i, \phi_i(t)}}(\gamma)| = |1 - \gamma(\phi_i(t))| |\widehat{f_i}(\gamma)| \geq \eta.
\]
This implies that $|1 - \gamma(\phi_i(t))| \geq \eta$ and $\gamma \in \Lambda_i := \{ \lambda \in \Gamma : |\widehat{f_i}(\lambda)| \geq \eta /2 \}$.

We have thus shown that 
\[
  B( \bigcup_{i=1}^k \Lambda_i \circ \phi_i; \eta) \subset B,  
\]
where $\Lambda_i \circ \phi_i: = \{ \gamma \circ \phi_i: \gamma \in \Lambda_i\} \subset \Gamma$. By  Parseval's formula, 
\[
\left( \frac{\eta}{2} \right)^2 |\Lambda_i| \leq \sum_{\lambda \in \Lambda_i} \left|\widehat{f_i}(\lambda)\right|^2 \leq 1.
\]
Therefore, $|\Lambda_i| \leq \frac{4}{\eta^2}$ and $|\bigcup_{i=1}^k \Lambda_i \circ \phi_i| \leq \frac{4k}{\eta^2}$.
\end{proof}

